\documentclass[11pt]{article}
\usepackage[a4paper,margin=25mm]{geometry}
\usepackage{amsmath,amssymb,amsthm,array,booktabs,longtable}
\usepackage[T1]{fontenc}
\usepackage{lmodern}
\usepackage[hidelinks,breaklinks]{hyperref}
\usepackage{xurl}
\newtheorem{theorem}{Theorem}[section]
\newtheorem{lemma}[theorem]{Lemma}
\newtheorem{corollary}[theorem]{Corollary}
\newcommand{\ZZ}{\mathbb Z}
\newcommand{\QQ}{\mathbb Q}
\newcommand{\RR}{\mathbb R}
\newcommand{\eps}{\varepsilon}
\title{Ovsienko--Tabachnikov's dual Somos integrality conjectures:\\
a two-weight extension of Hone's divisibility proof}
\author{Haobo Ma\\\small ChronoAI Pte Ltd, Singapore
\and Wenlin Zhang\\\small National University of Singapore\\\small ChronoAI Pte Ltd, Singapore}
\date{17 November 2026}
\begin{document}
\maketitle
\begin{abstract}
Extending Hone's divisibility proof (2007, Proposition~1)~\cite{Hone07}
and Ovsienko and Tabachnikov's sufficient conditions, we settle the
conjecture immediately after Corollary~5.2 in Section~5.4, printed
pages~10--11 of arXiv:1705.01623v1 (2017)~\cite{OT}, by classifying
integer shadows of both dual Somos deformations. For positive integers
$p,q$, integer weights $u,v$, and each fixed integer initial shadow tuple,
the recurrence
$A_{n+4}A_n=(1+u\eps)(A_{n+1}A_{n+3})^p+(1+v\eps)A_{n+2}^q$
has integer shadows at every index exactly when
$(q-2)u=(2p-2)v$. In the first conjectured family, $u=0,v=1$,
this is equivalent to $p=1$; if $p\ge2$, the first noninteger shadow
is always at index $9$. In the second family, $u=1,v=0$, the criterion
is $q=2$. More generally, the first failure and its reduced denominator
depend only on the real orbit and the integer obstruction, not on the
initial shadows. We also give the arbitrary-initial-shadow version of
the source's $q=0$ example and an exact certificate for its reported
sign change in the fourth homogeneous shadow.
\end{abstract}

\section{The recurrence and the result}
For $K=\ZZ,\QQ$, or $\RR$, put
$D(K)=K[\eps]/(\eps^2)$. Its elements are uniquely $a+b\eps$,
with multiplication
$(a+b\eps)(c+d\eps)=ac+(ad+bc)\eps$.
We call $a$ the real part and $b$ the shadow. In $D(\QQ)$ and
$D(\RR)$ an element is invertible exactly when $a\ne0$;
its inverse is $a^{-1}-ba^{-2}\eps$.
Fix $p,q\in\ZZ_{>0}$, $u,v\in\ZZ$, and
$c=(c_0,c_1,c_2,c_3)\in\ZZ^4$. Set $U=1+u\eps$ and $V=1+v\eps$.
\begin{equation}\label{eq:rec}
 A_{n+4}A_n=U(A_{n+1}A_{n+3})^p+VA_{n+2}^q,\quad
 A_j=1+c_j\eps\quad(0\le j\le3).
\end{equation}
All indices are nonnegative. Integer shadows mean $b_n\in\ZZ$ when
$A_n=a_n+b_n\eps$. We use indices $0,\ldots,3$ for the four leading
slots, as in the beginning of Section~5.4 of~\cite{OT}.

\begin{theorem}\label{thm:main}
The recurrence~\eqref{eq:rec} has a unique orbit in $D(\QQ)$ and a unique
orbit in $D(\RR)$; the latter is its coefficientwise image. For every
choice of the parameters above, all $a_n$ are positive integers and
each four consecutive $a_n$ are pairwise coprime. For each fixed $c$,
\begin{equation}\label{eq:iff}
 \bigl(b_n\in\ZZ\text{ for every }n\ge0\bigr)
 \quad\Longleftrightarrow\quad
 (q-2)u=(2p-2)v.
\end{equation}
\end{theorem}

Ovsienko and Tabachnikov proved the sufficient directions for the
two singly weighted families~\cite[Corollary~5.2]{OT}, using weighted
quivers and a Laurent phenomenon. Their preceding discussion also
gives integer shadows in the homogeneous case $u=v=0$.
For the first sufficient family, the earlier coefficient Laurent theorem
of Hone~\cite[Section~3, Proposition~1]{Hone07} already gives $p=1,u=0$
with arbitrary $v,c$, by specializing its coefficient $\beta$ to
$1+v\eps$ and its initial variables to $1+c_j\eps$.
The Laurent phenomenon~\cite{FZ} and mutation-periodic
quivers~\cite{FM} provide the broader setting. For the classical
exponents $p=1,q=2$, arbitrary dual coefficients and their Laurent
property follow from Fomin--Zelevinsky~\cite[Theorem~1.8]{FZ} by dual
specialization and occur explicitly in Hone~\cite[Section~1, Lemma~1.1]{Hone}.
Our proof adapts Hone's earlier coprimality-and-congruence
argument~\cite{Hone07} to the quotient of the integer dual-number ring.
Theorem~\ref{thm:main} supplies the necessary and sufficient condition
for all positive $p,q$, including simultaneous weights.
The tuple $c$ is fixed separately in each equivalence, not quantified
inside an all-initial-data integrality predicate.
The original first-family conjecture says ``is integer if and only if
$p=1$''~\cite[p.~11]{OT}. The second family is interpreted with the
exponent $q$ used in the proof and Example~5.3; the literal printed
formula is distinguished below.

\section{The local obstruction}
Write
\begin{equation}\label{eq:kappa}
 \kappa=(2p-2)v-(q-2)u.
\end{equation}
We work in $R=D(\ZZ)$. For any integer $w$ and any integer exponent $m$,
$(1+w\eps)^m=1+mw\eps$; negative powers use
$(1+w\eps)^{-1}=1-w\eps$.

\begin{lemma}\label{lem:local}
Let $A,B,C,D,E,F,G,H\in R$ satisfy
\begin{align*}
 EA&=U(BD)^p+VC^q,& FB&=U(CE)^p+VD^q,\\
 GC&=U(DF)^p+VE^q,& HD&=U(EG)^p+VF^q.
\end{align*}
Suppose the real part $e$ of $E$ is nonzero and coprime to the real
parts of $B,C,D$. Then
\begin{equation}\label{eq:local}
 E\mid U(FH)^p+VG^q\quad\Longleftrightarrow\quad e\mid\kappa.
\end{equation}
Divisibility on the left is in $R$ and on the right in $\ZZ$.
\end{lemma}
\begin{proof}
Pass to $S=R/(E)$. Coprimality implies that the images of $B,C,D$
are units. Indeed, if $b$ is the real part of $B$, choose integers
$r,s$ with $rb+se=1$. Then $rB+sE$ has real part $1$, so is a unit
in $R$. Its image is $rB$, showing that $B$ is a unit in $S$.
The same argument applies to $C,D$. The elements $U,V$ are already
units in $R$.

Omitting bars in $S$, the four equations become
\begin{equation}\label{eq:quotient}
 (BD)^p=-VU^{-1}C^q,\qquad
 F=VB^{-1}D^q,\qquad
 G=UC^{-1}(DF)^p,\qquad
 H=VD^{-1}F^q.
\end{equation}
Thus $F,G,H$ are also units. Ratios in the next calculation denote
multiplication by inverses of these units:
\begin{align*}
 \frac{VG^q}{U(FH)^p}
 &=U^{q-1}V^{1-p}C^{-q}D^{pq+p}F^{-p}\\
 &=U^{q-1}V^{1-2p}C^{-q}(BD)^p\\
 &=-U^{q-2}V^{2-2p}.
\end{align*}
The middle equality uses $BF=VD^q$, and the last uses the first
equation of~\eqref{eq:quotient}. Consequently,
\begin{equation}\label{eq:key}
 U(FH)^p+VG^q
 =U(FH)^p\bigl(1-U^{q-2}V^{2-2p}\bigr)
 =U(FH)^p\kappa\eps\quad\text{in }S.
\end{equation}
The factor $U(FH)^p$ is a unit, so the left side vanishes in $S$
exactly when $\kappa\eps$ does. Finally,
$E\mid\kappa\eps$ in $R$ is equivalent to $e\mid\kappa$ in $\ZZ$:
in an equality $E(r+s\eps)=\kappa\eps$, the real part gives $er=0$,
so $r=0$ and $es=\kappa$. Conversely $es=\kappa$ gives
$E(s\eps)=\kappa\eps$. This proves~\eqref{eq:local} without cancelling
$\eps$.
\end{proof}

\section{Integral orbits and the converse}
\begin{lemma}\label{lem:real}
The rational real-part recurrence
\begin{equation}\label{eq:real}
 a_{n+4}a_n=(a_{n+1}a_{n+3})^p+a_{n+2}^q,\qquad
 a_0=a_1=a_2=a_3=1
\end{equation}
is everywhere positive and satisfies $a_{n+1}>a_n$ for $n\ge3$.
Whenever a prefix is integral, each four consecutive terms in that
prefix are pairwise coprime.
\end{lemma}
\begin{proof}
Positive previous values give a positive next value, so division in
\eqref{eq:real} is always defined. Initially the values are nondecreasing
and at least $1$. If the values through $a_{n+3}$ have those properties,
then, since $p\ge1$,
$(a_{n+1}a_{n+3})^p\ge a_na_{n+3}$, and the remaining term in the
numerator is positive. Hence $a_{n+4}>a_{n+3}$, proving the claim.

For coprimality, suppose $a,b,c,d$ are four consecutive integral
terms, pairwise coprime, and the next term $e$ is integral. Then
$ae=(bd)^p+c^q$. If a prime divides both $e$ and $b$, it divides
$c^q$, contrary to $\gcd(b,c)=1$. The same argument handles $e,d$.
If a prime divides $e,c$, it divides $(bd)^p$, contrary to the
coprimality of $c$ with both $b$ and $d$. Thus $b,c,d,e$ are
pairwise coprime. The initial window consists of four ones.
\end{proof}

\begin{proof}[Proof of Theorem~\ref{thm:main}]
The real parts satisfy~\eqref{eq:real}, independently of $u,v,c$.
Lemma~\ref{lem:real} makes every preceding dual denominator a unit in
$D(\QQ)$ and in $D(\RR)$. Forward division therefore gives existence
and uniqueness in both rings. Induction identifies the real orbit
with the coefficientwise image of the rational orbit.

First assume $\kappa=0$. The four initial dual values are units of
$R=D(\ZZ)$, so division by them constructs $A_4,\ldots,A_7$ in $R$.
Inductively suppose all values through $A_{n+3}$ lie in $R$, where
$n\ge4$. Lemma~\ref{lem:real} supplies coprimality of all real
four-term windows in this prefix. Apply Lemma~\ref{lem:local} to
\[
 (A,B,C,D,E,F,G,H)
 =(A_{n-4},A_{n-3},A_{n-2},A_{n-1},A_n,A_{n+1},A_{n+2},A_{n+3}).
\]
The four required exchange equations have indices $n-4,\ldots,n-1$.
Since $a_n\mid0$, the next numerator is divisible by $A_n$ in $R$.
The quotient equals $A_{n+4}$ by uniqueness of division in $D(\QQ)$.
This proves integrality of the whole dual orbit for $\kappa=0$.
In particular, choosing $u=v=0$ and $c=0$ proves that the sequence
$a_n$ is integral. Since real parts are parameter-independent, it is
integral for every $u,v,c$; Lemma~\ref{lem:real} now gives all the
claimed coprimality statements.

Conversely, suppose every $b_n$ is an integer. The already proved
integrality of the real parts implies $A_n\in R$ for all $n$.
Applying Lemma~\ref{lem:local} at every $n\ge4$, and using the actual
next term as the quotient, gives $a_n\mid\kappa$ for every $n\ge4$.
The positive integers $a_n$ strictly increase from index $3$, so are
unbounded. They cannot all divide a fixed nonzero integer. Thus
$\kappa=0$, as required.
\end{proof}

\begin{corollary}\label{cor:failure}
Suppose $\kappa\ne0$, and let
\[
 m=\min\{n\ge4:a_n\nmid\kappa\}.
\]
The first noninteger shadow is $b_{m+4}$. Its denominator in lowest
terms is
\[
 \frac{a_m}{\gcd(a_m,\kappa)}.
\]
Both the index and the denominator are independent of $c$. In particular,
the failure index is at most $|\kappa|+7$. If $\kappa$ is odd, it is $8$,
and the denominator is $2$.
\end{corollary}
\begin{proof}
The sequence $a_n$ is unbounded, so $m$ exists. Starting with the
integral prefix through $A_7$, Lemma~\ref{lem:local} gives precisely
the asserted first failure. To determine its denominator, set
$E=A_m=e+e'\eps$, and let $X=x+x'\eps$ be its next numerator.
Its real quotient $t=x/e=a_{m+4}$ is integral. In $S=R/(E)$,
$X=(x'-te')\eps$, whereas~\eqref{eq:key} writes the same class as a
unit times $\kappa\eps$. The cyclic additive subgroup generated by
$\eps$ has order $e$, since $E\mid k\eps$ exactly when $e\mid k$.
Multiplication by a unit preserves additive order. Thus the order of
$(x'-te')\eps$ equals $e/\gcd(e,\kappa)$, which is exactly the reduced
denominator of $b_{m+4}=(x'-te')/e$. Finally $a_n\ge n-2$ for $n\ge4$;
at $n=|\kappa|+3$ it exceeds $|\kappa|$. The odd case uses $a_4=2$.
\end{proof}

\begin{corollary}\label{cor:weights}
Homogeneous shadows ($u=v=0$) are integral for all positive $p,q$.
For $p=1,q=2$, every integer pair $(u,v)$ is admissible. Otherwise put
$g=\gcd(2p-2,|q-2|)>0$. The admissible pairs are exactly
\[
 (u,v)=t\left(\frac{2p-2}{g},\frac{q-2}{g}\right),\qquad t\in\ZZ.
\]
\end{corollary}
\begin{proof}
Apply Theorem~\ref{thm:main}. In the last case the two displayed
coordinates are coprime, so $(q-2)u=(2p-2)v$ gives exactly these
integer multiples, including when one coordinate is zero.
\end{proof}

\section{The two conjectured families}
For fixed $p,q$, write $b^{u,v,c}$ for the shadow sequence and let
$h^{(j)}=b^{0,0,e_j}$, where $e_j$ is the $j$th standard initial tuple.
\begin{corollary}\label{cor:affine}
For all positive $p,q$,
\[
 b^{u,v,c}_n=u b^{1,0,0}_n+v b^{0,1,0}_n
              +\sum_{j=0}^3 c_jh^{(j)}_n.
\]
In the classical case $p=1,q=2$,
$\sum_{j=0}^3 h^{(j)}_n=a_n$ and
$\sum_{j=0}^3 jh^{(j)}_n=na_n$.
\end{corollary}
\begin{proof}
The $\eps$ coefficient of~\eqref{eq:rec} is a linear recurrence in
the shadows with forcing linear in $u,v$, and its leading coefficient
$a_n$ is positive. Superposition and the four initial conditions give
the first identity. For $p=1,q=2$, both $(a_n)$ and $(na_n)$ solve
the homogeneous shadow recurrence: the scalar factors in its three
products are respectively $2$, and $2n+4$. Their initial shadows
are $(1,1,1,1)$ and $(0,1,2,3)$, giving the remaining identities by
uniqueness. The first of these is the decomposition in
\cite[Section~2.1]{OT}; the affine decomposition is described in
\cite[Section~2.2]{OT}.
\end{proof}

\begin{corollary}\label{cor:first}
For $u=0,v=1$ and every fixed $c\in\ZZ^4$, all shadows are integers
if and only if $p=1$. If $p\ge2$, the first noninteger shadow is
$b_9$, for every positive $q$ and every $c$.
\end{corollary}
\begin{proof}
The equivalence follows from $\kappa=2p-2$.
All values through $A_7$ are integral because the first four
denominators are units of $R$. Since $a_4=2$ divides $2p-2$,
Lemma~\ref{lem:local} at $n=4$ gives $A_8\in R$.
But $a_5=2^p+1$, and at $n=5$ the same lemma gives
$A_9\in R$ exactly when $2^p+1\mid2p-2$. For $p\ge2$,
\[
 0<2p-2<2^p+1.
\]
For completeness, $2^p\ge2p$ for $p\ge1$: it holds at $p=1$,
and doubling $2^p\ge2p$ gives $2^{p+1}\ge4p\ge2(p+1)$.
Thus the divisibility fails. Since $a_9$ is an integer, its shadow
$b_9$ is not.
Its reduced denominator is
$(2^p+1)/\gcd(2^p+1,2p-2)$ by Corollary~\ref{cor:failure}.
\end{proof}

\begin{corollary}\label{cor:second}
For the second deformation
\[
 A_{n+4}A_n=(1+\eps)(A_{n+1}A_{n+3})^p+A_{n+2}^q,
\]
and each fixed $c\in\ZZ^4$, all shadows are integers if and only if
$q=2$.
\end{corollary}
\begin{proof}
Set $u=1,v=0$ in Theorem~\ref{thm:main}; then $\kappa=2-q$.
In particular every odd positive $q$ gives its first noninteger shadow
at index $8$, with denominator $2$, for every $p,c$.
\end{proof}
Here the second deformation is parameterized as in Example~5.3
of~\cite{OT}. The printed recurrence (18) has exponent $1$ on its
last factor and says ``is integer'', whereas its proof specifies
$q=2$. The literal exponent-$1$ assertion fails: $p=1,c=0$ gives
$b_8=159/2$, already recorded in Example~5.3(b). Replacing the last
factor by $A_{n+2}^2$ gives the sufficient condition proved there.
This is a correction of the displayed formula, not a counterexample
to the stated conjecture with its intended parameter $q$.

\begin{corollary}\label{cor:zero}
For $p=1,q=0,u=1,v=0$, the first noninteger shadow is $b_9$ for
every $c\in\ZZ^4$, and its reduced denominator is $3$.
\end{corollary}
\begin{proof}
Here the real prefix is $(1,1,1,1,2,3,4,9,14,19)$.
Substitution into the shadow recurrence gives the following exact
affine expressions $b_n=d_n+L_nc$, where $L_nc$ is the dot product.
\begin{center}\small
\begin{tabular}{@{}crrrrr@{}}
\toprule
$n$ & $d_n$ & \multicolumn{4}{c}{$L_n$}\\
\midrule
4 & 1 & $-2$ & 1 & 0 & 1\\
5 & 3 & $-2$ & $-2$ & 2 & 1\\
6 & 6 & $-2$ & $-2$ & $-2$ & 4\\
7 & 24 & $-12$ & 0 & $-4$ & 3\\
8 & 56 & $-13$ & $-16$ & 3 & 2\\
9 & $307/3$ & $-14$ & $-18$ & $-18$ & 15\\
\bottomrule
\end{tabular}
\end{center}
These rows follow successively from
$a_nb_{n+4}=a_{n+1}b_{n+3}+a_{n+3}b_{n+1}
+a_{n+1}a_{n+3}-a_{n+4}b_n$, so each equality can be checked
without the positive-$q$ theorem. The first five rows are integral
and the final row has fractional part $1/3$ for every integer $c$.
The specialization $c=0$ is Example~5.3(a) of~\cite{OT}.
\end{proof}

\section{A finite sign obstruction}
The homogeneous classical shadow with initial tuple $(0,0,0,1)$
has
\[
 b_{25}=-545740823162385126878041<0,
\]
where indices start at zero. This confirms the earlier negative-sign
observation in~\cite[Section~2.1, p.~3]{OT}, whose indices start at one.
It also gives an explicit counterexample to the later conjecture
in~\cite[Example~3.8]{Hone} that the fourth homogeneous shadow is
``positive for all $n\ge2$'': Hone's $y^{(\mathrm{iv})}_{24}=b_{25}$.
Indeed his initial indices $-1,0,1,2$ correspond to our $0,1,2,3$,
with the same four initial shadows and homogeneous recurrence.

The file \texttt{shadow-certificate.csv} supplies the exact integer
pairs $(a_n,b_n)$ for $0\le n\le25$. Each successive row satisfies
\[
 a_na_{n+4}=a_{n+1}a_{n+3}+a_{n+2}^2,\qquad
 a_nb_{n+4}=a_{n+1}b_{n+3}+2a_{n+2}b_{n+2}
 +a_{n+3}b_{n+1}-a_{n+4}b_n.
\]
Together with the initial data, these identities certify the displayed
negative term. In this prefix, $b_n>0$ for $3\le n\le24$.
This finite sign obstruction does not assert the eventual sign of any
homogeneous shadow or settle the two forced-shadow positivity and
growth conjectures in~\cite[Section~2.2]{OT}.

\clearpage
\section{Coverage and supporting calculations}
\small
\begin{longtable}{@{}>{\raggedright\arraybackslash}p{0.42\textwidth}>{\raggedright\arraybackslash}p{0.53\textwidth}@{}}
\toprule
Source statement & Mathematical coverage\\
\midrule
\endfirsthead
\toprule
Source statement & Mathematical coverage (continued)\\
\midrule
\endhead
\cite{OT}, Section 5.4, p.~10, Corollary 5.2(i), $p=1$
 & Reproved; follows by dual substitution from Hone (2007,
 Proposition~1); explicit dual proof Ovsienko--Tabachnikov (2017).\\[3pt]
\cite{OT}, Section 5.4, p.~11, conjecture after the proof, first deformation
 & Proved for every initial tuple; Corollary~\ref{cor:first} also locates the first failure.\\[3pt]
Same conjecture, second deformation, and Example 5.3(b,c), $q=1,3$
 & Proved for $q>0$ in Corollary~\ref{cor:second}; odd $q$ fails first at $8$.
 The $q=2$ sufficient direction is reproved; prior proof
 Ovsienko--Tabachnikov (2017).\\[3pt]
\cite{OT}, p.~10, literal displayed (18), last exponent $1$
 & Refuted-and-corrected: $p=1,c=0,b_8=159/2$;
 corrected last exponent $2$, as in its proof.\\[3pt]
\cite{OT}, Section 1.2 (order four), Section 2.1 integrality and basis sum,
and Section 2.2 affine decomposition
 & Reproved; homogeneous integrality follows from Fomin--Zelevinsky
 (2001 preprint, Example~3.3); the decompositions are in
 Ovsienko--Tabachnikov (2017). Corollaries~\ref{cor:weights} and~\ref{cor:affine}.\\[3pt]
\cite{OT}, p.~11, Example 5.3(a), $p=1,q=0$
 & Prior example Ovsienko--Tabachnikov (2017); extended to every
 initial tuple in Corollary~\ref{cor:zero}.\\[3pt]
\cite{OT}, Section 2.1, fourth homogeneous shadow's reported negative terms
 & Prior sign observation Ovsienko--Tabachnikov (2017);
 exact finite certificate through zero-based index $25$.\\[3pt]
\cite{Hone}, Lemma 1.1, arbitrary classical dual coefficients
 & Reproved at real coefficients $1,1$; dual specialization of
 Fomin--Zelevinsky (2001, Theorem~1.8), explicit dual treatment by
 Hone (2021 preprint; 2023 journal). Earlier singly weighted examples:
 Ovsienko~\cite{Ovs}.\\[3pt]
\cite{Hone}, Example 3.8, fourth homogeneous shadow positivity conjecture
 & Refuted (explicit counterexample): $y^{(\mathrm{iv})}_{24}<0$.
 Negative-sign observation already in Ovsienko--Tabachnikov (2017).\\[3pt]
Simultaneous-weight lattice; exact first-failure index and denominator
 & Not posed in~\cite{OT}; proved in
 Corollaries~\ref{cor:weights} and~\ref{cor:failure}.\\[3pt]
\cite{OT}, Section 2.1, other homogeneous-shadow signs;
 Section 2.2, forced-shadow positivity and faster growth
 & Not covered. Hone~\cite[Example~3.8]{Hone} proves positivity for a
 different homogeneous basis, not these two forced shadows.\\[3pt]
\cite{OT}, Sections 1.1--1.2, other orders; Sections 5.2--5.3,
 Lucas and order-three examples (OEIS A001871, A005246)
 & Not covered; the proof uses four-term windows. For dual Somos-5,
 see Harrow--Hone~\cite{HH}.\\[3pt]
\cite{OT}, Section 5.5, asymptotics and complete integrability question
 & Not covered by the divisibility proof. Classical analytic solutions
 are given in Hone~\cite{Hone}; no general-exponent conclusion is inferred.\\
\bottomrule
\end{longtable}
\normalsize

The accompanying \texttt{verify.py} uses exact rational arithmetic and
two implementations of the dual recurrence. Its finite coverage is
$1\le p\le3$, $1\le q\le4$, $-2\le u,v\le2$, indices $0\le n\le9$,
and $83$ initial tuples: all of $\{-1,0,1\}^4$, together with
$(4,-7,11,-13)$ and $(-9,0,2,17)$. It checks coefficient agreement,
positive integral real parts, coprime windows, the local criterion
whenever its preceding eight terms are integral, the reduced-denominator
formula, and the first failure in Corollary~\ref{cor:first}.
It also checks the affine $q=0$ table on these initial tuples and all
rows of the homogeneous-shadow certificate using both implementations.
These finite calculations supplement the
infinite written proof; their exact counts are in \texttt{verification.json}.
For example, $p=q=1,u=1,v=0,c=0$ gives $b_8=159/2$, agreeing with
Example~5.3(b) of~\cite{OT} after indexing its leading slots from zero.

The pinned Lean coverage is the first-family statement
$u=0,v=1$ for every positive $p,q$ and every integer initial tuple,
including existence, uniqueness, positivity, and the rational-to-real
dual-number correspondence. Its declaration in the following fixed source is
\begin{center}\small
\path{D5.S1.Recurrence.Invariants.SomosFirstDualIntegrality.result}
\end{center}
\begin{center}\footnotesize
\url{https://github.com/the-omega-institute/trureturing/blob/256870690452e5629b1260ae620f2d647417e857/D5/S1/Recurrence/Invariants/SomosFirstDualIntegrality.lean}
\end{center}
Here positivity means the real part $a_n>0$, not positivity of a
shadow. The definitions \texttt{IsOrbit} and \texttt{IsRealOrbit} use
the four slots $0,1,2,3$, real initial values $1$, arbitrary integer
shadows, and the first deformation $u=0,v=1$.
The two-weight classification, second-family converse, exact
first-failure denominator, and $q=0$ assertion have the written proofs
above. The sign counterexample has the finite certificate just stated.
These scopes are distinct from that declaration's coverage.

\small
\begin{thebibliography}{9}
\bibitem{OT} V.~Ovsienko and S.~Tabachnikov,
\emph{Dual numbers, weighted quivers, and extended Somos and Gale--Robinson sequences},
Algebr. Represent. Theory \textbf{21} (2018), 1119--1132.
\href{https://doi.org/10.1007/s10468-018-9779-3}{doi:10.1007/s10468-018-9779-3}.
Posing version: \href{https://arxiv.org/abs/1705.01623v1}{arXiv:1705.01623v1} (2017).
\bibitem{FZ} S.~Fomin and A.~Zelevinsky,
\emph{The Laurent phenomenon}, Adv. Appl. Math. \textbf{28} (2002), 119--144.
\href{https://doi.org/10.1006/aama.2001.0770}{doi:10.1006/aama.2001.0770}.
\bibitem{FM} A.~P.~Fordy and R.~J.~Marsh,
\emph{Cluster mutation-periodic quivers and associated Laurent sequences},
J. Algebraic Combin. \textbf{34} (2011), 19--66.
\href{https://doi.org/10.1007/s10801-010-0262-4}{doi:10.1007/s10801-010-0262-4}.
\bibitem{Hone} A.~N.~W.~Hone,
\emph{Casting light on shadow Somos sequences},
Glasg. Math. J. \textbf{65} (2023), suppl.~S1, S87--S101.
\href{https://doi.org/10.1017/S0017089522000167}{doi:10.1017/S0017089522000167}.
First preprint: \href{https://arxiv.org/abs/2111.10905v1}{arXiv:2111.10905v1} (2021).
\bibitem{Hone07} A.~N.~W.~Hone,
\emph{Laurent polynomials and superintegrable maps},
SIGMA \textbf{3} (2007), 022, 18~pp.
\href{https://doi.org/10.3842/SIGMA.2007.022}{doi:10.3842/SIGMA.2007.022}.
\bibitem{Ovs} V.~Ovsienko,
\emph{A step towards cluster superalgebras},
\href{https://arxiv.org/abs/1503.01894v7}{arXiv:1503.01894v7} (2016),
Section~7, Corollary~7.1 and the following second extension.
\bibitem{HH} J.~W.~E.~Harrow and A.~N.~W.~Hone,
\emph{Casting more light in the shadows: dual Somos-5 sequences},
\href{https://arxiv.org/abs/2409.00406}{arXiv:2409.00406} (2024).
\end{thebibliography}
\medskip\noindent\small
\textcopyright\ 2026 Haobo Ma and Wenlin Zhang.
Licensed under Creative Commons Attribution 4.0 International (CC BY 4.0).
\end{document}
